% Project Description — <<TITLE>>
% Call: <<CALL>> (read <<CALL_READ_DATE>>). Governing rules: PAPPG <<PAPPG>> II.D.2.d.
% Limits from CALL.md: <<PD_PAGES>> pages total; Results from Prior NSF Support
% at most <<PRIOR_PAGES>> of them; no URLs (name repositories in words; link
% them in the Data Management Plan). Required headings: "Results from Prior
% NSF Support" with "Intellectual Merit" and "Broader Impacts" inside it;
% "Broader Impacts" as a section heading on its own line.
% Style (using-pdpowers/references/PD-POLICY.md): third person ("the PI",
% "this proposal"); short declarative sentences; the PI's own results in
% prose with \cite; main definitions in `definition`, main conjectures in
% `conjecture`; every project rests on a cited PI result and says in prose
% what will exist afterwards; no run-in labels, no fixed-count step lists.
% Labels: \item\label{item:<name>} on every enumerate item that is referred
% to, cited as \eqref{item:<name>}; never "item 3" in prose.
\documentclass[11pt]{amsart}
\usepackage{amsmath,amssymb,amsfonts,amsthm}
\usepackage[letterpaper,margin=1in]{geometry}
\usepackage{enumitem}
\usepackage{graphicx}
\usepackage{color}
\usepackage{mathtools}
\usepackage{booktabs}
\pagenumbering{arabic}
\newcommand{\taylor}[1]{{\color{red} \sf $\spadesuit\spadesuit\spadesuit$ Taylor: [#1]}}
\theoremstyle{plain}
\newtheorem{theorem}{Theorem}[section]
\newtheorem{conjecture}[theorem]{Conjecture}
\theoremstyle{definition}
\newtheorem{definition}[theorem]{Definition}
\newtheorem{problem}{Problem}
\newtheorem{project}{Project}
\newtheorem{question}{Question}
\numberwithin{equation}{section}

\begin{document}
\title[<<SHORT_TITLE>>]{<<TITLE>>}
\maketitle
\vspace{-1cm}

%% ------------------------------------------------------------------
\section{Overview}\label{S:overview}
%% ------------------------------------------------------------------
% About one page. In the order the panel reads it:
% 1. The subject and the question in two sentences.
% 2. What is already in hand (each result with its classical anchor).
% 3. The method sentence: who drives, what the tools do; the
%    solicitation's language in at most one short quotation.
% 4. The list of projects, one clause each, with \ref to their subsections.
% 5. What the proposal produces (the solicitation's measures of success:
%    publications, data sets, tools, students trained).

%% ------------------------------------------------------------------
\section{Background}\label{S:background}
%% ------------------------------------------------------------------
\subsection{Introduction to the subject}
% Definitions the panel needs, built by analogy with a classical subject.
% Each main definition in \begin{definition}...\end{definition}; the
% classical statement in `theorem` (at most two) and the conjecture that
% generalizes it in `conjecture`, with attribution and date.

\subsection{History and the state of the field}
% Dated list of partial progress; who proved what; what is open now.

\subsection{Previous work of the PI}
% The PI's results the projects rest on, with citations. Published,
% submitted with statement intact, or PI-verified only; everything else
% is "agents report" or belongs in REVIEW.md.

%% ------------------------------------------------------------------
\section{Results from Prior NSF Support}\label{S:prior}
%% ------------------------------------------------------------------
% PAPPG II.D.2.d(iii). At most <<PRIOR_PAGES>> pages. One award: the most
% closely related, current or ended within five years.
\subsection*{<<PRIOR_AWARD_NUMBER>>, \emph{<<PRIOR_AWARD_TITLE>>}, <<PRIOR_AWARD_AMOUNT>>, <<PRIOR_AWARD_PERIOD>>}
% (a) number, amount, period; (b) title — above.
\subsubsection*{Intellectual Merit}
% (c) numbered products, each with \cite; results in prose.
\begin{enumerate}[leftmargin=*,itemsep=1pt]
\item\label{item:prior-<<PRIOR_ITEM_1_SLUG>>}
\end{enumerate}
\subsubsection*{Broader Impacts}
% (c) continued: what was done (students, software released, lectures).
\subsubsection*{Publications and products}
% (d) list of publications (full citations here or in References Cited),
% or the sentence "No publications were produced under this award."
% (e) research products and where they live (software, data, video).
\subsubsection*{Relation to the proposed work}
% (f) one or two sentences: which items each project starts from, as
%     \eqref{item:prior-<<PRIOR_ITEM_1_SLUG>>}.

%% ------------------------------------------------------------------
\section{Method and Standards}\label{S:method}
%% ------------------------------------------------------------------
% How the work is done and what counts as a result. Define every term the
% panel has not seen before (e.g. a dump, a saturated manuscript, fidelity
% of a formalization). State the AI usage categories and how usage is
% recorded. State the practices; do not argue for them or quote the
% solicitation at length. The reader draws the conclusion that the tools
% are "scientifically justified rather than incidental".
\subsection{Standards every project meets}
\begin{enumerate}[leftmargin=*,itemsep=1pt]
\item
\end{enumerate}
\subsection{Tools and records}

%% ------------------------------------------------------------------
\section{Proposed Research}\label{S:proposed}
%% ------------------------------------------------------------------
% One \subsection per project, titled by its object or theorem, never an
% activity. Plain prose, each paragraph opening with its topic sentence; no
% run-in labels ("Built on:", "What is new.", "Plan."), no "Step 1" items.
% Paragraphs in this order (pd-project): foundation (the PI result(s) this
% rests on, cited, and the external theorems used); objective (one sentence
% naming the theorem, object, database or tool that will exist); preliminary
% results at verified strength (a small table of real records if data); plan
% as prose or \subsubsection headings, each part naming its technique (a
% numbered list only for a proof program); feasibility (what has run, at
% what size, where the cost sits); dated status lines naming what is missing;
% what the project yields if its question resolves either way; deliverables
% and compute in prose. Name prior constructions of the same
% kind and the senior people whose results are used; cite them fairly. A
% conjecture the project attacks goes in `conjecture`, attributed.
\subsection{<<PROJECT_1_TITLE>>}\label{S:project1}

\subsection{<<PROJECT_2_TITLE>>}\label{S:project2}

%% ------------------------------------------------------------------
\section{Intellectual Merit}\label{S:merit}
%% ------------------------------------------------------------------
% One paragraph: why these problems, why now, why this PI.

%% ------------------------------------------------------------------
\section{Broader Impacts}\label{S:impact}
%% ------------------------------------------------------------------
% Heading must read exactly "Broader Impacts" (PAPPG). One paragraph per
% item (textbook; open-source tools and data; students; seminars and
% visitors; lectures online), each opening with the done thing, dated, then
% its continuation. Plain paragraphs; no \textbf{...} run-in labels.


\newpage\label{S:bibstart}
% Bibliography here is for citation resolution only; build.sh trims it from
% description.pdf. References Cited is a separate document.
\bibliographystyle{amsalpha}
\bibliography{references}
\end{document}
